
\SGARefTarget{sga3:iv:section:4}\section*{4. Topologies and Sheaves}
\addcontentsline{toc}{section}{4. Topologies and Sheaves}

The notion of a sieve, and the presentation of the notion of a topology
adopted here (\SGARefLink{sga3:iv:definition:4:2:1}{4.2.1}), which is more intrinsic and in many respects more
convenient than the presentation by covering families in \SGARefLink{sga3:iv:bib:ma}{[MA]}, are due to
J.~Giraud \SGARefLink{sga3:iv:bib:as}{[AS]}.

\SGARefTarget{sga3:iv:subsection:4:1}\subsection*{4.1. Sieves}
\addcontentsline{toc}{subsection}{\SGARefLink{sga3:iv:subsection:4:1}{4.1}. Sieves}

\medskip
\noindent\SGARefTarget{sga3:iv:definition:4:1:1}\textbf{Definition 4.1.1.}
\textit{A \emph{sieve} of the category \(\mathscr C\) is a subfunctor
\(C\) of the final functor}
\[
  e:\mathscr C^{\circ}\longrightarrow\mathbf{Set}.
\]

To every sieve \(C\) of \(\mathscr C\) one associates the set \(E(C)\) of
objects \(X\) of \(\mathscr C\) such that \(C(X)\ne\varnothing\), that is,
such that the structural morphism \(X\to e\) factors through \(C\). Thus
\[
\tag{+}\label{sga3:iv:definition:4:1:1:eq:unnamed:01}
\left\{
\begin{aligned}
 X\in E(C)&\quad\Longleftrightarrow\quad C(X)=e(X)=\{\varnothing\},\\
 X\notin E(C)&\quad\Longleftrightarrow\quad C(X)=\varnothing.
\end{aligned}
\right.
\]
The set \(E=E(C)\) has \SGARefLink{sga3:iv:definition:4:1:1:eq:unnamed:02}{the following property}:
\[
\tag{++}\label{sga3:iv:definition:4:1:1:eq:unnamed:02}
 X\in E\ \text{and}\ \operatorname{Hom}(Y,X)\ne\varnothing
 \quad\Longrightarrow\quad Y\in E.
\]
(If the set \(\operatorname{Ob}\mathscr C\) is endowed with its natural
preorder, with \(Y\) dominating \(X\) when there is an arrow from \(Y\) to
\(X\), then the sets \(E\) satisfying \SGARefLink{sga3:iv:definition:4:1:1:eq:unnamed:02}{(++)} are precisely the subsets of
\(\operatorname{Ob}\mathscr C\) that contain every object dominating one of
their elements.)

Conversely, if \(E\) is a subset of \(\operatorname{Ob}\mathscr C\) having
property \SGARefLink{sga3:iv:definition:4:1:1:eq:unnamed:02}{(++)}, then \(E\) can be written uniquely in the form \(E(C)\), with
\(C\) defined by the formulas \SGARefLink{sga3:iv:definition:4:1:1:eq:unnamed:01}{(+)}. Thus sieves of \(\mathscr C\) correspond
bijectively to subsets of \(\operatorname{Ob}\mathscr C\) satisfying \SGARefLink{sga3:iv:definition:4:1:1:eq:unnamed:02}{(++)}.
By abuse of language, we shall sometimes call the set \(E(C)\) itself a
sieve of \(\mathscr C\).

Let \(C\) and \(C'\) be two sieves of \(\mathscr C\).\footnote{Editor's
note: The following sentence has been expanded.} Since both are subfunctors
of the final functor \(e\), it is equivalent to say that \(C\) dominates
\(C'\) (for the domination relation in
\(\operatorname{Ob}\widehat{\mathscr C}\)), that \(C\) is a subfunctor of
\(C'\), or that \(E(C)\subset E(C')\). In this case we say that \(C\) is
\emph{finer} than \(C'\). This is an order relation on the set of sieves of
\(\mathscr C\). Moreover,
\[
 E(C)\cap E(C')=E(C\times C'),
\]
so the set of sieves of \(\mathscr C\) is filtered for the order relation
``is finer than.''

Every subset \(E\) of \(\operatorname{Ob}\widehat{\mathscr C}\), and in
particular every subset of \(\operatorname{Ob}\mathscr C\), defines a sieve
\(C(E)\): the set of \(X\in\operatorname{Ob}\mathscr C\) such that
\(F(X)\ne\varnothing\) for at least one \(F\in E\) satisfies \SGARefLink{sga3:iv:definition:4:1:1:eq:unnamed:02}{(++)}, and hence
defines the desired sieve.

This sieve may also be defined as the image of the family of morphisms
\(\{F\to e\mid F\in E\}\), in the sense of \SGARefLink{sga3:iv:definition:4:1:2}{the following definition}.

\medskip
\noindent\SGARefTarget{sga3:iv:definition:4:1:2}\textbf{Definition 4.1.2.}
\textit{Let \(\{F_i\to F\}\) be a family of morphisms of
\(\widehat{\mathscr C}\) with common target \(F\). The \emph{image} of this
family is the subfunctor of \(F\) defined by}
\[
 S\longmapsto\bigcup_i\operatorname{Im}F_i(S)\subset F(S).
\]

\medskip
\noindent\SGARefTarget{sga3:iv:proposition:4:1:3}\textbf{Proposition 4.1.3.}
\textit{Formation of the image commutes with base change. In the preceding
notation let \(I\) be the image of the family \(\{F_i\to F\}\). For every
morphism \(G\to F\) in \(\widehat{\mathscr C}\), the image of the family}
\[
 \{F_i\times_FG\longrightarrow G\}
\]
\textit{is the subfunctor \(I\times_FG\) of \(G\).}

\medskip
\noindent\SGARefTarget{sga3:iv:definition:4:1:4:0}\textbf{Definition 4.1.4.0.}
\textit{Let \(C\) be a sieve of \(\mathscr C\). If \(E\) is a subset of
\(\operatorname{Ob}\mathscr C\) such that \(C(E)=C\), then \(E\) is called
a \emph{base} of \(C\). Every sieve \(C\) has a base, for example
\(E(C)\).}

We now describe \(\operatorname{Hom}(C,F)\), where \(C\) is a sieve of
\(\mathscr C\) and \(F\in\operatorname{Ob}\widehat{\mathscr C}\), by means
of a base \(\{S_i\}\) of \(C\). For every pair \((i,j)\) there is a
diagram in \(\widehat{\mathscr C}\):
\[
\begin{array}{ccc}
 S_i\times S_j&\longrightarrow&S_i\\
 \big\downarrow&&\big\downarrow\\
 S_j&\longrightarrow&e.
\end{array}
\]
It gives a diagram of sets
\[
 \Gamma(F)=\operatorname{Hom}(e,F)
 \xrightarrow{\sigma}
 \prod_i\operatorname{Hom}(S_i,F)
 \mathrel{\substack{\xrightarrow{\tau_1}\\[-0.5ex]
                     \xrightarrow[\tau_2]{}}}
 \prod_{i,j}\operatorname{Hom}(S_i\times S_j,F),
\]
with \(\tau_1\sigma=\tau_2\sigma\). Hence there is a morphism
\[
 \operatorname{Hom}(e,F)\longrightarrow
 \operatorname{Ker}\!\left(
 \prod_i\operatorname{Hom}(S_i,F)
 \rightrightarrows
 \prod_{i,j}\operatorname{Hom}(S_i\times S_j,F)
 \right).
\]
One checks immediately:

\medskip
\noindent\SGARefTarget{sga3:iv:proposition:4:1:4}\textbf{Proposition 4.1.4.}
\textit{There is an isomorphism, functorial in \(F\),}
\[
 \operatorname{Hom}(C,F)\xrightarrow{\sim}
 \operatorname{Ker}\!\left(
 \prod_i\operatorname{Hom}(S_i,F)
 \rightrightarrows
 \prod_{i,j}\operatorname{Hom}(S_i\times S_j,F)
 \right),
\]
\textit{for which the diagram}
\[
\begin{array}{ccc}
 \operatorname{Hom}(e,F)&\longrightarrow&
 \operatorname{Ker}\!\left(
 \prod_i\operatorname{Hom}(S_i,F)\rightrightarrows
 \prod_{i,j}\operatorname{Hom}(S_i\times S_j,F)\right)\\
 \Big\Vert&&\overset{\sim}{\Big\uparrow}\\
 \operatorname{Hom}(e,F)&\longrightarrow&\operatorname{Hom}(C,F)
\end{array}
\]
\textit{is commutative; the bottom arrow is induced by the canonical
morphism \(C\to e\).}

\medskip
\noindent\SGARefTarget{sga3:iv:corollary:4:1:5}\textbf{Corollary 4.1.5.}
\textit{Suppose that the fiber products \(S_i\times S_j\) are representable,
for example that the \(S_i\) are squarable. Then, for every
\(F\in\operatorname{Ob}\widehat{\mathscr C}\), there is an isomorphism}
\[
 \operatorname{Hom}(C,F)\xrightarrow{\sim}
 \operatorname{Ker}\!\left(
 \prod_iF(S_i)\rightrightarrows
 \prod_{i,j}F(S_i\times S_j)\right).
\]

\medskip
\noindent\SGARefTarget{sga3:iv:remark:4:1:6}\textbf{Remark 4.1.6.}
Let \(R\) be a sieve of \(\mathscr C\), let \(\overline R\) denote the full
subcategory of \(\mathscr C\) whose object set is \(E(R)\), and let
\[
 i_R:\overline R\longrightarrow\mathscr C
\]
be the inclusion functor. There is an isomorphism, functorial in
\(F\in\operatorname{Ob}\widehat{\mathscr C}\),
\[
 \operatorname{Hom}(R,F)\xrightarrow{\sim}\Gamma(F\circ i_R),
\]
for which the diagram
\[
\begin{array}{ccc}
 \operatorname{Hom}(e,F)&\longrightarrow&\operatorname{Hom}(R,F)\\
 \overset{\sim}{\Big\downarrow}&&
 \overset{\sim}{\Big\downarrow}\\
 \Gamma(F)&\longrightarrow&\Gamma(F\circ i_R)
\end{array}
\]
is commutative; the bottom arrow is induced by \(i_R\).

\medskip
\noindent\SGARefTarget{sga3:iv:definition:4:1:7}\textbf{Definition 4.1.7.}
\textit{Let \(\mathscr C\) be a category. A \emph{sieve of the object}
\(S\) of \(\mathscr C\) is a sieve of the category
\(\mathscr C_{/S}\).}

A sieve of \(S\) is therefore a subobject of \(S\) in
\(\widehat{\mathscr C}\). It corresponds canonically to a subset of
\(\operatorname{Ob}\mathscr C_{/S}\) that contains the source of every arrow
whose target it contains. By abuse of language, such a subset will also be
called a sieve of \(S\).

\SGARefTarget{sga3:iv:subsection:4:2}\subsection*{4.2. Topologies: Definitions}
\addcontentsline{toc}{subsection}{\SGARefLink{sga3:iv:subsection:4:2}{4.2}. Topologies: Definitions}

\medskip
\noindent\SGARefTarget{sga3:iv:definition:4:2:1}\textbf{Definition 4.2.1.}
\textit{A \emph{topology} on a category \(\mathscr C\) consists, for every
object \(S\) of \(\mathscr C\), of a set \(J(S)\) of sieves of \(S\), called
\emph{covering sieves} or \emph{refinements} of \(S\), satisfying the
following axioms:}
\begin{enumerate}
 \item[(T 1)]\SGARefTarget{sga3:iv:definition:4:2:1:item:t:1:01} \textit{For every refinement \(R\) of \(S\) and every
 morphism \(T\to S\), the sieve \(R\times_ST\) of \(T\) is covering
 (``stability under base change'').}
 \item[(T 2)]\SGARefTarget{sga3:iv:definition:4:2:1:item:t:2:01} \textit{If \(R,C\) are two sieves of \(S\), if \(R\) is
 covering, and if for every \(T\in\operatorname{Ob}\mathscr C\) and every
 morphism \(T\to R\), the sieve \(C\times_ST\) of \(T\) is covering, then
 \(C\) is a refinement of \(S\).}
 \item[(T 3)]\SGARefTarget{sga3:iv:definition:4:2:1:item:t:3:01} \textit{If \(C\supset R\) are two sieves of \(S\) and
 \(R\) is covering, then \(C\) is covering.}
 \item[(T 4)]\SGARefTarget{sga3:iv:definition:4:2:1:item:t:4:01} \textit{For every \(S\), \(S\) is a refinement of \(S\).}
\end{enumerate}

These axioms may be reformulated as follows. Suppose that a topology
\(S\mapsto J(S)\) on \(\mathscr C\) is given. For every
\(F\in\operatorname{Ob}\widehat{\mathscr C}\), let \(J(F)\) be the set of
subfunctors \(R\) of \(F\) such that, for every morphism \(T\to F\) in
\(\widehat{\mathscr C}\) with \(T\) representable, the sieve
\(R\times_FT\) of \(T\) is covering. By \SGARefLink{sga3:iv:definition:4:2:1:item:t:1:01}{(T~1)}, this notation agrees with
the preceding one. An \(R\in J(F)\) is again called a \emph{refinement} of
\(F\). The preceding axioms immediately imply:
\begin{enumerate}
 \item[\((T'0)\)]\SGARefTarget{sga3:iv:definition:4:2:1:item:tprime0:01} If \(F\supset G\) are objects of
 \(\widehat{\mathscr C}\), and if
 \(G\times_FS\in J(S)\) for every \(S\in\operatorname{Ob}\mathscr C\) and
 every morphism \(S\to F\), then \(G\in J(F)\).
 \item[\((T'1)\)]\SGARefTarget{sga3:iv:definition:4:2:1:item:tprime1:01} If \(G\in J(F)\) and \(H\to F\) is a morphism in
 \(\widehat{\mathscr C}\), then \(G\times_FH\in J(H)\).
 \item[\((T'2)\)]\SGARefTarget{sga3:iv:definition:4:2:1:item:tprime2:01} If \(F\supset G\supset H\), if \(G\in J(F)\),
 and if \(H\in J(G)\), then \(H\in J(F)\).
 \item[\((T'3)\)]\SGARefTarget{sga3:iv:definition:4:2:1:item:tprime3:01} If \(F\supset G\supset H\) and \(H\in J(F)\),
 then \(G\in J(F)\).
 \item[\((T'4)\)]\SGARefTarget{sga3:iv:definition:4:2:1:item:tprime4:01} For every
 \(F\in\operatorname{Ob}\widehat{\mathscr C}\), one has \(F\in J(F)\).
\end{enumerate}
Conversely, if for every \(F\in\operatorname{Ob}\widehat{\mathscr C}\)
one is given a set \(J(F)\) of subobjects of \(F\) satisfying
\SGARefLink{sga3:iv:definition:4:2:1:item:tprime0:01}{\(T'0\)}--\SGARefLink{sga3:iv:definition:4:2:1:item:tprime4:01}{\(T'4\)}, then \(S\mapsto J(S)\) defines a topology on
\(\mathscr C\), and these two constructions are inverse to each other.

Properties \SGARefLink{sga3:iv:definition:4:2:1:item:tprime1:01}{\(T'1\)}, \SGARefLink{sga3:iv:definition:4:2:1:item:tprime2:01}{\(T'2\)}, and \SGARefLink{sga3:iv:definition:4:2:1:item:tprime3:01}{\(T'3\)} imply the further
property:
\begin{enumerate}
 \item[\((T'5)\)]\SGARefTarget{sga3:iv:definition:4:2:1:item:tprime5:01} If \(G\) and \(H\) are two subobjects of \(F\) and
 \(G,H\in J(F)\), then \(G\cap H\in J(F)\).
\end{enumerate}
Thus \(J(F)\), ordered by \(\supset\), is filtered. This observation will
be useful below.

\medskip
\noindent\SGARefTarget{sga3:iv:para:4:2:2}\textbf{4.2.2.}
The topology defined by \(J\) is said to be \emph{finer} than the topology
defined by \(J'\) if
\(J(S)\supset J'(S)\) for every
\(S\in\operatorname{Ob}\mathscr C\). Equivalently,
\(J(F)\supset J'(F)\) for every
\(F\in\operatorname{Ob}\widehat{\mathscr C}\).

Every set of topologies on \(\mathscr C\) has a greatest lower bound. More
precisely, let \(I\) be an index set and, for each \(i\in I\), let
\(S\mapsto J_i(S)\) be a topology on \(\mathscr C\). Then
\[
 J(S)=\bigcap_{i\in I}J_i(S)
\]
defines a topology, and it is the greatest lower bound of the given set.

In particular, suppose that for every
\(S\in\operatorname{Ob}\mathscr C\) a set \(E(S)\) of sieves of \(S\) is
given. The \emph{topology generated by these sets} is the least fine
topology for which every element of \(E(S)\) is a refinement of \(S\).

\medskip
\noindent\SGARefTarget{sga3:iv:definition:4:2:3}\textbf{Definition 4.2.3.}
\textit{Let \(\{F_i\to F\}\) be a family of morphisms in
\(\widehat{\mathscr C}\), and let \(G\subset F\) be its image (\SGARefLink{sga3:iv:definition:4:1:2}{4.1.2}). The
family is \emph{covering} if \(G\in J(F)\). A morphism is covering if the
one-member family consisting of that morphism is covering.}

\SGARefLink{sga3:iv:definition:4:2:3}{This definition} applies in particular to an inclusion: a sieve \(C\) of
\(S\) is covering if and only if the canonical morphism \(C\to S\) is
covering.

Axioms \SGARefLink{sga3:iv:definition:4:2:1:item:tprime0:01}{\(T'0\)}--\SGARefLink{sga3:iv:definition:4:2:1:item:tprime5:01}{\(T'5\)} give the following properties of covering
families:
\begin{enumerate}
 \item[(C 0)]\SGARefTarget{sga3:iv:definition:4:2:3:item:c:0:01} If \(\{F_i\to F\}\) is a family in
 \(\widehat{\mathscr C}\) such that every representable base change
 \(\{F_i\times_FS\to S\}\) is covering, then the original family is
 covering.
 \item[(C 1)]\SGARefTarget{sga3:iv:definition:4:2:3:item:c:1:01} If \(\{F_i\to F\}\) is covering and \(G\to F\) is any
 morphism, then \(\{F_i\times_FG\to G\}\) is covering (``stability under
 base change'').
 \item[(C 2)]\SGARefTarget{sga3:iv:definition:4:2:3:item:c:2:01} If \(\{F_i\to F\}\) is covering and, for every \(i\),
 \(\{F_{ij}\to F_i\}\) is covering, then the composite family
 \(\{F_{ij}\to F\}\) is covering (``stability under composition'').
 \item[(C 3)]\SGARefTarget{sga3:iv:definition:4:2:3:item:c:3:01} If \(\{G_j\to F\}\) is covering and \(\{F_i\to F\}\) is
 a family such that for every \(j\) there is an \(i\) for which
 \(G_j\to F\) factors through \(F_i\to F\), then \(\{F_i\to F\}\) is
 covering (``saturation'').
 \item[(C 4)]\SGARefTarget{sga3:iv:definition:4:2:3:item:c:4:01} Every one-member family consisting of an isomorphism is
 covering.
\end{enumerate}
Notice that \SGARefLink{sga3:iv:definition:4:2:3:item:c:2:01}{(C~2)} and \SGARefLink{sga3:iv:definition:4:2:3:item:c:3:01}{(C~3)} also imply:
\begin{enumerate}
 \item[(C 5)]\SGARefTarget{sga3:iv:definition:4:2:3:item:c:5:01} If \(\{F_i\to F\}\) is a family for which there exists a
 covering family \(\{G_j\to F\}\) such that
 \(\{F_i\times_FG_j\to G_j\}\) is covering for every \(j\), then
 \(\{F_i\to F\}\) is covering (``a locally covering family is
 covering'').
\end{enumerate}

\medskip
\noindent\SGARefTarget{sga3:iv:para:4:2:4}\textbf{4.2.4.}
Conversely, let \(\mathscr C\) be a category having fiber products, and
suppose that for every \(S\in\operatorname{Ob}\mathscr C\) one is given a
set of families of morphisms in \(\mathscr C\) with target \(S\), called
covering families, satisfying \SGARefLink{sga3:iv:definition:4:2:3:item:c:1:01}{(C~1)}--\SGARefLink{sga3:iv:definition:4:2:3:item:c:4:01}{(C~4)}, and hence also \SGARefLink{sga3:iv:definition:4:2:3:item:c:5:01}{(C~5)}. Let
\(J(S)\) be the set of sieves of \(S\) having a covering base (or,
equivalently by \SGARefLink{sga3:iv:definition:4:2:3:item:c:3:01}{(C~3)}, all of whose bases are covering). Then
\(S\mapsto J(S)\) defines a topology on \(\mathscr C\), and the two
preceding constructions are inverse to one another.

In applications it is inconvenient to consider all covering families, since
one often has a fairly simple description of a ``sufficient'' number of
them. This leads to the following definitions.

\medskip
\noindent\SGARefTarget{sga3:iv:definition:4:2:5:0}\textbf{Definition 4.2.5.0.}\footnote{Editor's note: \SGARefLink{sga3:iv:definition:4:2:5:0}{This
definition}, which appeared after \SGARefLink{sga3:iv:definition:4:2:5}{4.2.5} in the original, has been placed here
because it will be used in \SGARefLink{sga3:iv:proposition:6:2:1}{6.2.1} in a slightly more general setting than
that of \SGARefLink{sga3:iv:definition:4:2:5}{4.2.5}.}
\textit{Let \(\mathscr C\) be a category, and suppose given, for every
\(S\in\operatorname{Ob}\mathscr C\), a set \(P(S)\) of families of
morphisms in \(\mathscr C\) with target \(S\). The \emph{topology generated
by \(P\)} is the least fine topology for which the given families are
covering.}

\medskip
\noindent\SGARefTarget{sga3:iv:definition:4:2:5}\textbf{Definition 4.2.5.}
\textit{A \emph{pretopology} on a category \(\mathscr C\) consists, for
every \(S\in\operatorname{Ob}\mathscr C\), of a set \(R(S)\) of
families \(\{S_i\to S\}\), called \emph{covering for the pretopology},
satisfying:}
\begin{enumerate}
 \item[(P 1)]\SGARefTarget{sga3:iv:definition:4:2:5:item:p:1:01} \textit{For every \(\{S_i\to S\}\in R(S)\) and
 every morphism \(T\to S\), the fiber products \(S_i\times_ST\) exist and
 \(\{S_i\times_ST\to T\}\in R(T)\).}
 \item[(P 2)]\SGARefTarget{sga3:iv:definition:4:2:5:item:p:2:01} \textit{If \(\{S_i\to S\}\in R(S)\) and, for each
 \(i\), \(\{T_{ij}\to S_i\}\in R(S_i)\), then the composite
 family \(\{T_{ij}\to S\}\) belongs to \(R(S)\).}
 \item[(P 3)]\SGARefTarget{sga3:iv:definition:4:2:5:item:p:3:01} \textit{Every one-member family consisting of an isomorphism
 is covering.}
\end{enumerate}

\medskip
\noindent\SGARefTarget{sga3:iv:proposition:4:2:6}\textbf{Proposition 4.2.6.}
\textit{For each \(S\), let \(J(S)\) be the set of sieves of \(S\) that
are covering for the topology generated by the pretopology \(R\).
Let \(J_R(S)\) be the subset of \(J(S)\) consisting of the
sieves defined by the families in \(R(S)\). Then
\(J_R(S)\) is cofinal in \(J(S)\): every refinement of \(S\)
contains a sieve defined by a family in \(R(S)\).}

For every \(S\), let \(J'(S)\) be the set of sieves of \(S\) containing a
sieve in \(J_R(S)\). Clearly \(J'(S)\subset J(S)\). To prove
\(J(S)=J'(S)\), it suffices to prove that the \(J'(S)\) define a topology,
that is, satisfy \SGARefLink{sga3:iv:definition:4:2:1:item:t:1:01}{(T~1)}--\SGARefLink{sga3:iv:definition:4:2:1:item:t:4:01}{(T~4)}. Axioms \SGARefLink{sga3:iv:definition:4:2:1:item:t:1:01}{(T~1)}, \SGARefLink{sga3:iv:definition:4:2:1:item:t:3:01}{(T~3)}, and \SGARefLink{sga3:iv:definition:4:2:1:item:t:4:01}{(T~4)} are immediate;
only \SGARefLink{sga3:iv:definition:4:2:1:item:t:2:01}{(T~2)} remains.

Let \(U\in J'(S)\), let \(C\) be a sieve of \(S\), and suppose that
\(C\times_ST\in J'(T)\) for every \(T\to U\). We must prove
\(C\in J'(S)\). By definition, \(U\) contains a refinement \(U'\) defined
by a family \(\{S_i\to S\}\in R(S)\). Since \SGARefLink{sga3:iv:definition:4:2:1:item:t:3:01}{(T~3)} has already been
verified, it is enough to prove \(U'\cap C\in J'(S)\); hence we may suppose
\(U=U'\). By hypothesis,
\(C\times_SS_i\in J'(S_i)\) for every \(i\). Thus for each \(i\) there is
a covering family \(\{T_{ij}\to S_i\}\in R(S_i)\) such that
\(T_{ij}\to S_i\) factors through \(C\times_SS_i\to S_i\). Consequently
\(T_{ij}\to S\) factors through \(C\to S\). Hence \(C\) contains the sieve
defined by the composite family \(\{T_{ij}\to S\}\), and \SGARefLink{sga3:iv:definition:4:2:5:item:p:2:01}{(P~2)} completes
the proof.

Axioms \SGARefLink{sga3:iv:definition:4:2:5:item:p:1:01}{(P~1)}--\SGARefLink{sga3:iv:definition:4:2:5:item:p:3:01}{(P~3)} are those of \SGARefLink{sga3:iv:bib:ma}{[MA]}. Because pretopologies are so useful
in practice, we shall interpret every important result in terms of a
pretopology defining the topology under consideration.

\medskip
\noindent\SGARefTarget{sga3:iv:remark:4:2:7}\textbf{Remark 4.2.7.}
One may introduce a slightly more general notion: for every \(S\), one
specifies a set of covering families satisfying \SGARefLink{sga3:iv:definition:4:2:5:item:p:1:01}{(P~1)}, \SGARefLink{sga3:iv:definition:4:2:5:item:p:3:01}{(P~3)}, and the
conclusion of Proposition~\SGARefLink{sga3:iv:proposition:4:2:6}{4.2.6}. This occurs in particular when the given
families satisfy \SGARefLink{sga3:iv:definition:4:2:5:item:p:1:01}{(P~1)}, \SGARefLink{sga3:iv:definition:4:2:5:item:p:3:01}{(P~3)}, and \SGARefLink{sga3:iv:definition:4:2:3:item:c:5:01}{(C~5)}. See \SGARefLink{sga3:iv:bib:d}{[D]}.

\medskip
\noindent\SGARefTarget{sga3:iv:definition:4:2:8}\textbf{Definition 4.2.8.}
\textit{Let \(\mathscr C\) be equipped with a topology, let \(S\) be an
object of \(\mathscr C\), and let \(P(S')\) be a condition on an argument
\(S'\in\operatorname{Ob}\mathscr C_{/S}\). Suppose that
\(\operatorname{Hom}(S'',S')\ne\varnothing\) implies
\(P(S')\Rightarrow P(S'')\). The condition \(P\) is said to hold
\emph{locally on \(S\)} for the topology under consideration if the
following equivalent conditions hold:}
\begin{enumerate}
 \item[(i)]\SGARefTarget{sga3:iv:definition:4:2:8:item:i:01} \textit{The set of \(S'\to S\) for which \(P(S')\) holds is a
 refinement of \(S\).}
 \item[(ii)]\SGARefTarget{sga3:iv:definition:4:2:8:item:ii:01} \textit{There is a refinement of \(S\) such that \(P(S')\)
 holds for every \(S'\) in that refinement.}
 \item[(iii)]\SGARefTarget{sga3:iv:definition:4:2:8:item:iii:01} \textit{If the topology is defined by a pretopology, there
 is a covering family for that pretopology such that \(P(S')\) holds for
 every \(S'\) in the family.}
\end{enumerate}

\medskip
\noindent\SGARefTarget{sga3:iv:example:4:2:9}\textbf{Example 4.2.9.}
\textit{Let \(f:X\to Y\) be an \(S\)-morphism. We say that \(f\) is
\emph{locally an isomorphism} if there is a covering family
\(\{S_i\to S\}\) such that \(f\times_SS_i\) is an isomorphism for every
\(i\). Equivalently, there is a refinement \(R\) of \(S\) such that
\(X(T)\to Y(T)\) is an isomorphism for every \(T\to R\).}

Many other examples of this ``local'' language will appear below.

\SGARefTarget{sga3:iv:subsection:4:3}\subsection*{4.3. Presheaves, Sheaves, and the Sheaf Associated to a Presheaf}
\addcontentsline{toc}{subsection}{\SGARefLink{sga3:iv:subsection:4:3}{4.3}. Presheaves, Sheaves, and the Sheaf Associated to a Presheaf}

\medskip
\noindent\SGARefTarget{sga3:iv:definition:4:3:1}\textbf{Definition 4.3.1.}
\textit{Let \(\mathscr C\) be a category. A \emph{presheaf of sets} on
\(\mathscr C\) is a contravariant functor from \(\mathscr C\) to the
category of sets. The category}
\[
 \widehat{\mathscr C}
 =\operatorname{Hom}(\mathscr C^{\circ},\mathbf{Set})
\]
\textit{is called the \emph{category of presheaves} on \(\mathscr C\). If
\(\mathscr C\) is equipped with a topology, a presheaf \(P\) is
\emph{separated} (respectively, is a \emph{sheaf}) if, for every
\(S\in\operatorname{Ob}\mathscr C\) and every \(R\in J(S)\), the canonical
map}
\[
\tag{+}\label{sga3:iv:definition:4:3:1:eq:unnamed:01}
 P(S)=\operatorname{Hom}(S,P)\longrightarrow\operatorname{Hom}(R,P)
\]
\textit{is injective (respectively, bijective). The full subcategory of
\(\widehat{\mathscr C}\) whose objects are the sheaves is called the
\emph{category of sheaves} and is denoted by
\(\widetilde{\mathscr C}\).}

\medskip
\noindent\SGARefTarget{sga3:iv:proposition:4:3:2}\textbf{Proposition 4.3.2.}
\textit{Let \(P\) be a separated presheaf (respectively, a sheaf). For
every \(H\in\operatorname{Ob}\widehat{\mathscr C}\) and every
\(R\in J(H)\), the canonical map}
\[
\tag{+}\label{sga3:iv:proposition:4:3:2:eq:unnamed:01}
 \operatorname{Hom}(H,P)\longrightarrow\operatorname{Hom}(R,P)
\]
\textit{is injective (respectively, bijective).}

Let \(P\) be separated, let \(H\) be a presheaf, let \(R\in J(H)\), and
let \(u,v:H\to P\) satisfy \(uj=vj\). For every \(f:S\to H\), with
\(S\in\operatorname{Ob}\mathscr C\), the sieve \(R\times_HS\) is a
refinement of \(S\), and \(ufj_S=vfj_S\). This is expressed by the diagram
\[
\begin{array}{ccccc}
 R\times_HS&\xrightarrow{j_S}&S&&\\
 \overset{f_R}{\big\downarrow}&&
 \overset{f}{\big\downarrow}&&\\
 R&\xrightarrow{j}&H&
 \mathrel{\substack{\xrightarrow{u}\\[-0.5ex]\xrightarrow[v]{}}}&P.
\end{array}
\]
Since \(P\) is separated, it follows that \(uf=vf\). This holds for every
representable \(S\), so \(u=v\).

Now suppose that \(P\) is a sheaf. Let \(g:R\to P\); we show that it
factors through \(H\). For every \(f:S\to H\), with
\(S\in\operatorname{Ob}\mathscr C\), the morphism
\(g\circ f_R:R\times_HS\to P\) factors uniquely through \(S\), thereby
defining a morphism \(h:S\to P\). By uniqueness, \(h\) is plainly functorial
in \(f\):
\[
\begin{tikzcd}[column sep=4.5em,row sep=3.1em]
 R\times_HS
   \arrow[r,"f_R"]
   \arrow[d,"j_S"']
 &
 R
   \arrow[r,"g"]
   \arrow[d,"j"']
 &
 P
 \\
 S
   \arrow[r,"f"']
   \arrow[rru,bend right=38,"h"']
 &
 H
   \arrow[ru,dashed]
 &
 {}
\end{tikzcd}
\]
Thus, for every \(S\), we have defined a map \(H(S)\to P(S)\) functorial
in \(S\), and hence a morphism \(H\to P\) with the required properties.

\medskip
\noindent\SGARefTarget{sga3:iv:corollary:4:3:2:1}\textbf{Corollary 4.3.2.1.}\footnote{Editor's note: \SGARefLink{sga3:iv:corollary:4:3:2:1}{This
corollary} has been added.}
\textit{Let \(R\) and \(F\) be two sheaves. If \(R\) is a refinement of
\(F\), then \(R=F\).}

Indeed, suppose that \(R\) is a refinement of \(F\), and write
\(j:R\hookrightarrow F\) for the inclusion. By \SGARefLink{sga3:iv:proposition:4:3:2}{4.3.2},
\(\operatorname{Hom}(F,R)=\operatorname{End}(R)\); hence there is a
\(\pi:F\to R\) such that \(\pi j=\operatorname{id}_R\). Likewise,
\(\operatorname{End}(F)=\operatorname{Hom}(R,F)\), and
\(j\pi j=j\) implies \(j\pi=\operatorname{id}_F\). Thus \(j\) is an
isomorphism.

\medskip
\noindent\SGARefTarget{sga3:iv:proposition:4:3:3}\textbf{Proposition 4.3.3 (\SGARefLink{sga3:iv:bib:as}{[AS]}, \SGARefLink{sga3:iv:definition:1:3}{1.3}).}
\textit{Let \(\mathscr C\) be a category and \(P\) a presheaf on
\(\mathscr C\). For every \(S\in\operatorname{Ob}\mathscr C\), let
\(J(S)\) be the set of sieves \(R\) of \(S\) such that, for every
\(T\to S\), the map}
\[
\tag{+}\label{sga3:iv:proposition:4:3:3:eq:unnamed:01}
 \operatorname{Hom}(T,P)\longrightarrow
 \operatorname{Hom}(R\times_ST,P)
\]
\textit{is injective (respectively, bijective). Then the \(J(S)\) define a
topology on \(\mathscr C\); that is, they satisfy \SGARefLink{sga3:iv:definition:4:2:1:item:t:1:01}{(T~1)}--\SGARefLink{sga3:iv:definition:4:2:1:item:t:4:01}{(T~4)}.}

\medskip
\noindent\SGARefTarget{sga3:iv:corollary:4:3:4}\textbf{Corollary 4.3.4.}
\textit{For every \(S\in\operatorname{Ob}\mathscr C\), let \(K(S)\) be a
family of sieves satisfying \SGARefLink{sga3:iv:definition:4:2:1:item:t:1:01}{(T~1)}, and let \(P\) be a presheaf on
\(\mathscr C\). Then \(P\) is separated (respectively, a sheaf) for the
topology generated by the \(K(S)\) if and only if, for every
\(S\in\operatorname{Ob}\mathscr C\) and every \(R\in K(S)\), the canonical
map}
\[
\tag{+}\label{sga3:iv:corollary:4:3:4:eq:unnamed:01}
 \operatorname{Hom}(S,P)\longrightarrow\operatorname{Hom}(R,P)
\]
\textit{is injective (respectively, bijective).}

\medskip
\noindent\SGARefTarget{sga3:iv:corollary:4:3:5}\textbf{Corollary 4.3.5.}
\textit{For each \(S\in\operatorname{Ob}\mathscr C\), let
\(R(S)\) be a set of families of morphisms in \(\mathscr C\) with
target \(S\), satisfying \SGARefLink{sga3:iv:definition:4:2:5:item:p:1:01}{(P~1)}, for example a pretopology. Let \(P\) be a
presheaf on \(\mathscr C\). Then \(P\) is separated (respectively, a sheaf)
for the topology generated by \(R\) if and only if, for every
\(S\in\operatorname{Ob}\mathscr C\) and every family
\(\{S_i\to S\}\in R(S)\), the map}
\[
 P(S)\longrightarrow\prod_iP(S_i)
\]
\textit{is injective (respectively, the diagram}
\[
 P(S)\longrightarrow\prod_iP(S_i)
 \rightrightarrows\prod_{i,j}P(S_i\times_SS_j)
\]
\textit{is exact).}

\medskip
\noindent\SGARefTarget{sga3:iv:definition:4:3:6}\textbf{Definition 4.3.6.}
\textit{The \emph{canonical topology} on a category \(\mathscr C\) is the
finest topology for which every representable functor is a sheaf.}

\medskip
\noindent\SGARefTarget{sga3:iv:corollary:4:3:7}\textbf{Corollary 4.3.7.}
\textit{A sieve \(R\) of \(S\) is a refinement for the canonical topology
if and only if, for every morphism \(T\to S\) in \(\mathscr C\) and every
\(X\in\operatorname{Ob}\mathscr C\), the canonical map}
\[
 \operatorname{Hom}(T,X)\longrightarrow
 \operatorname{Hom}(R\times_ST,X)
\]
\textit{is bijective.}

\medskip
\noindent\SGARefTarget{sga3:iv:definition:4:3:8}\textbf{Definition 4.3.8.}
\textit{A sieve covering for the canonical topology is called a
\emph{universal effective epimorphic sieve}.}

\medskip
\noindent\SGARefTarget{sga3:iv:corollary:4:3:9}\textbf{Corollary 4.3.9.}
\textit{A universal effective epimorphic family defines a universal
effective epimorphic sieve. Conversely, every squarable family defining a
universal effective epimorphic sieve is universal effective epimorphic.}

Return now to a category \(\mathscr C\) equipped with an arbitrary topology
and to the construction of the sheaf associated to a presheaf \(P\). Let
\(S\) be an object of \(\mathscr C\). If \(R\supset R'\) are two
refinements of \(S\), there is a diagram
\[
\begin{array}{ccc}
 \operatorname{Hom}(S,P)&\longrightarrow&\operatorname{Hom}(R,P)\\
 &\searrow&\big\downarrow\\
 &&\operatorname{Hom}(R',P).
\end{array}
\]
As already observed, the ordered set \(J(S)\) is filtered. Since
\(S\in J(S)\), there is an evident morphism
\[
 \operatorname{Hom}(S,P)\longrightarrow
 \varinjlim_{R\in J(S)}\operatorname{Hom}(R,P).
\]

\medskip
\noindent\SGARefTarget{sga3:iv:definition:4:3:10:0}\textbf{Definition 4.3.10.0.}
Set
\[
 \check H^0(S,P)=
 \varinjlim_{R\in J(S)}\operatorname{Hom}(R,P).
\]
One checks that \(\check H^0(S,P)\) depends functorially on \(S\), and hence
defines a functor \(LP\) by
\[
\tag{++}\label{sga3:iv:definition:4:3:10:0:eq:unnamed:01}
 \operatorname{Hom}(S,LP)=\check H^0(S,P)
 =\varinjlim_{R\in J(S)}\operatorname{Hom}(R,P).
\]
By construction there are morphisms
\[
 \ell_P:P\longrightarrow LP,
 \qquad
 z_R:\operatorname{Hom}(R,P)\longrightarrow\operatorname{Hom}(S,LP).
\]

\medskip
\noindent\SGARefTarget{sga3:iv:lemma:4:3:10}\textbf{Lemma 4.3.10.}
\begin{enumerate}
 \item[(i)]\SGARefTarget{sga3:iv:lemma:4:3:10:item:i:01} \textit{For every refinement \(R\) of \(S\) and every
 \(u:R\to P\), the diagram}
 \[
 \begin{array}{ccc}
 P&\xrightarrow{\ell_P}&LP\\
 \overset{u}{\Big\uparrow}&&
 \overset{z_R(u)}{\Big\uparrow}\\
 R&\xrightarrow{i_R}&S
 \end{array}
 \]
 \textit{is commutative.}
 \item[(ii)]\SGARefTarget{sga3:iv:lemma:4:3:10:item:ii:01} \textit{For every morphism \(v:S\to LP\), there exist a
 refinement \(R\) of \(S\) and a morphism \(u:R\to P\) such that
 \(v=z_R(u)\).}
 \item[(iii)]\SGARefTarget{sga3:iv:lemma:4:3:10:item:iii:01} \textit{Let \(Q\) be a functor and let \(u,v:Q\to P\)
 satisfy \(\ell_Pu=\ell_Pv\). Then the kernel of the pair \((u,v)\) is a
 refinement of \(Q\).}
 \item[(iv)]\SGARefTarget{sga3:iv:lemma:4:3:10:item:iv:01} \textit{Let \(u:R\to P\) and \(u':R'\to P\). The equality
 \(z_R(u)=z_{R'}(u')\) holds if and only if there is a refinement
 \(R''\subset R\cap R'\) of \(S\) on which \(u\) and \(u'\) coincide.}
\end{enumerate}

\noindent\textit{Proof.}
For \SGARefLink{sga3:iv:lemma:4:3:10:item:i:01}{(i)}, we must prove \(z_R(u)i_R=\ell_Pu\). It is enough to compare the
composites of these two morphisms with every morphism \(g:T\to R\), where
\(T\) is representable. Put \(f=i_Rg\) and
\(R'=R\times_ST\). We have
\[
\begin{array}{ccccc}
 P&\xrightarrow{\ell_P}&LP&&\\
 \overset{u}{\Big\uparrow}&&
 \overset{z_R(u)}{\Big\uparrow}&&\\
 R&\xrightarrow{i_R}&S&&\\
 \overset{p}{\Big\uparrow}&\overset{g}{\nwarrow}&
 \overset{f}{\Big\uparrow}&&\\
 R'&\xrightarrow[\sim]{i_{R'}}&T.&&
\end{array}
\]
By \SGARefLink{sga3:iv:definition:4:3:10:0}{the definition} of \(\ell_P\),
\(\ell_Pug=z_R(u)f\); this is the special case of \SGARefLink{sga3:iv:lemma:4:3:10:item:i:01}{the assertion} in which
\(i_R\) is an isomorphism. But \(z_R(u)f=z_R(u)i_Rg\), proving \SGARefLink{sga3:iv:lemma:4:3:10:item:i:01}{(i)}.

Assertions \SGARefLink{sga3:iv:lemma:4:3:10:item:ii:01}{(ii)} and \SGARefLink{sga3:iv:lemma:4:3:10:item:iv:01}{(iv)} merely express \SGARefLink{sga3:iv:definition:4:3:10:0}{the definition} of
\(\operatorname{Hom}(S,LP)\) as a direct limit.

For \SGARefLink{sga3:iv:lemma:4:3:10:item:iii:01}{(iii)}, let \(K\) be the kernel of \((u,v)\). For every morphism
\(f:S\to Q\), with \(S\) representable, \(K\times_QS\) is a subfunctor of
the kernel of \(uf,vf:S\to P\). By \SGARefLink{sga3:iv:definition:4:2:1:item:tprime0:01}{\(T'0\)}, it is therefore enough to
prove \SGARefLink{sga3:iv:lemma:4:3:10:item:iii:01}{the assertion} when \(Q=S\) is representable. In that case, \SGARefLink{sga3:iv:lemma:4:3:10:item:ii:01}{(ii)} and
\SGARefLink{sga3:iv:lemma:4:3:10:item:iv:01}{(iv)} show that \(K\) contains a refinement of \(S\), and hence is itself a
refinement of \(S\).

One finally checks that \(P\mapsto LP\) defines a functor
\[
 L:\widehat{\mathscr C}\longrightarrow\widehat{\mathscr C},
\]
and that \(P\mapsto\ell_P\) defines a morphism of functors
\[
 \ell:\operatorname{Id}_{\widehat{\mathscr C}}\longrightarrow L.
\]

We now state the essential result.

\medskip
\noindent\SGARefTarget{sga3:iv:proposition:4:3:11}\textbf{Proposition 4.3.11.}
\begin{enumerate}
 \item[(i)]\SGARefTarget{sga3:iv:proposition:4:3:11:item:i:01} \textit{For every presheaf \(P\), the presheaf \(LP\) is
 separated and \(\ell_P:P\to LP\) is covering (\SGARefLink{sga3:iv:definition:4:2:3}{4.2.3}).}
 \item[(ii)]\SGARefTarget{sga3:iv:proposition:4:3:11:item:ii:01} \textit{If \(P\) is a sheaf, then \(P\to LP\) is an
 isomorphism.}
 \item[(iii)]\SGARefTarget{sga3:iv:proposition:4:3:11:item:iii:01} \textit{For every presheaf \(P\) and every separated
 presheaf (respectively, sheaf) \(F\), the map}
 \[
  \operatorname{Hom}(\ell_P,F):
  \operatorname{Hom}(LP,F)\longrightarrow\operatorname{Hom}(P,F)
 \]
 \textit{is injective (respectively, bijective).}
 \item[(iv)]\SGARefTarget{sga3:iv:proposition:4:3:11:item:iv:01} \textit{If \(P\) is separated, then
 \(\ell_P:P\to LP\) is a covering monomorphism (so \(P\) is a refinement
 of \(LP\)), and \(LP\) is a sheaf.}
\end{enumerate}

\noindent\textit{Proof.}
\SGARefLink{sga3:iv:proposition:4:3:11:item:i:01}{(i)} First, \(P\to LP\) is covering. Indeed, for every \(v:S\to LP\),
Lemma~\SGARefLink{sga3:iv:lemma:4:3:10:item:i:01}{4.3.10(i)},\SGARefLink{sga3:iv:lemma:4:3:10:item:ii:01}{(ii)} supplies a refinement \(R\) of \(S\) and a commutative
diagram (the map from \(R\) to the fiber product is induced by
\(v':R\to P\) and \(i_R:R\to S\))
\[
\begin{tikzcd}[column sep=4.2em,row sep=2.8em]
 {}
 &
 P
   \arrow[r,"\ell_P"]
 &
 LP
 \\
 {}
 &
 P\times_{LP}S
   \arrow[u]
   \arrow[r]
 &
 S
   \arrow[u,"v"']
 \\
 R
   \arrow[ru]
   \arrow[ruu,bend left=10,"v'"]
   \arrow[rru,bend right=8,"i_R"']
 &
 {}
 &
 {}
\end{tikzcd}
\]
so \(P\times_{LP}S\to S\) is covering. By \SGARefLink{sga3:iv:definition:4:2:3:item:c:0:01}{(C~0)}, \(P\to LP\) is
covering.

Now suppose that two morphisms \(v_1,v_2:S\to LP\) induce the same
morphism on a refinement \(R\) of \(S\). There are refinements \(R_i\),
\(i=1,2\), and morphisms \(u_i:R_i\to P\) with
\(z_{R_i}(u_i)=v_i\). Replacing \(R\) by a smaller refinement, we may
suppose \(R_1=R_2=R\). The commutative diagram of Lemma~\SGARefLink{sga3:iv:lemma:4:3:10:item:i:01}{4.3.10(i)} then
gives \(\ell_Pu_1=\ell_Pu_2\). By \SGARefLink{sga3:iv:lemma:4:3:10:item:iii:01}{4.3.10(iii)}, \(u_1\) and \(u_2\)
coincide on a refinement of \(R\), hence on a refinement of \(S\); by
\SGARefLink{sga3:iv:lemma:4:3:10:item:iv:01}{4.3.10(iv)}, \(z_R(u_1)=z_R(u_2)\). Thus \(LP\) is separated.

\SGARefLink{sga3:iv:proposition:4:3:11:item:ii:01}{(ii)} This is clear: if \(P\) is a sheaf, then
\(\operatorname{Hom}(S,P)\to\operatorname{Hom}(R,P)\) is already an
isomorphism for every refinement \(R\) of \(S\).

\SGARefLink{sga3:iv:proposition:4:3:11:item:iii:01}{(iii)} Let \(u,v:LP\to F\) satisfy \(u\ell_P=v\ell_P\). To show \(u=v\),
it is enough to prove \(uf=vf\) for every \(f:S\to LP\) with \(S\)
representable. There are a refinement \(R\) of \(S\) and a morphism
\(g:R\to P\) with \(f=z_R(g)\). The morphisms \(uf\) and \(vf\) agree on
\(R\), where both equal \(u\ell_Pg=v\ell_Pg\). If \(F\) is separated,
they therefore agree on \(S\), proving injectivity. If \(F\) is a sheaf,
the commutative diagram
\[
\begin{tikzcd}[column sep=large,row sep=large]
 P \arrow[r,"\ell_P"] \arrow[d] &
 LP \arrow[d] \arrow[dl]\\
 F \arrow[r,"\sim"] & LF
\end{tikzcd}
\]
shows that \(\operatorname{Hom}(\ell_P,F)\) is also surjective.

\SGARefLink{sga3:iv:proposition:4:3:11:item:iv:01}{(iv)} Suppose that \(P\) is separated. To show that \(P\to LP\) is a
monomorphism, it is enough to consider \(u,v:S\to P\), with \(S\)
representable, satisfying \(\ell_Pu=\ell_Pv\). By \SGARefLink{sga3:iv:lemma:4:3:10:item:iii:01}{4.3.10(iii)}, \(u\) and
\(v\) agree on a refinement of \(S\), and therefore agree because \(P\) is
separated. Thus \(P\to LP\) is a monomorphism; since it is covering by
\SGARefLink{sga3:iv:proposition:4:3:11:item:i:01}{(i)}, \(P\) is a refinement of \(LP\).

It remains to show that \(LP\) is a sheaf. We already know from \SGARefLink{sga3:iv:proposition:4:3:11:item:i:01}{(i)} that it
is separated. Let \(S\in\operatorname{Ob}\mathscr C\), let \(R\) be a
refinement of \(S\), and let \(h:R\to LP\). We must find \(u:S\to LP\)
with \(ui_R=h\). Since \(P\) is a refinement of \(LP\),
\(R'=P\times_{LP}R\) is a refinement of \(R\), and hence of \(S\). Put
\(u=z_{R'}(h')\); the situation is
\[
\begin{array}{ccccc}
 P&\xrightarrow{\ell_P}&LP&&\\
 \overset{h'}{\Big\uparrow}&&
 \overset{h}{\Big\uparrow}&\overset{u}{\nwarrow}&\\
 R'&\xrightarrow{j}&R&\xrightarrow{i_R}&S.
\end{array}
\]
Here the diagonal arrow is \(u:S\to LP\).
We have \(ui_{R'}=\ell_Ph'=hj\), hence \(ui_Rj=hj\). Since \(R'\) is a
refinement of \(R\) and \(LP\) is separated, Proposition~\SGARefLink{sga3:iv:proposition:4:3:2}{4.3.2} gives
\(ui_R=h\).

\medskip
\noindent\SGARefTarget{sga3:iv:corollary:4:3:12}\textbf{Corollary 4.3.12.}\footnote{Editor's note: The statement
and proof of \SGARefLink{sga3:iv:corollary:4:3:12}{this corollary} have been expanded.}
\textit{Let \(F\) be a sheaf and let \(R\) be a subobject of \(F\) in
\(\widehat{\mathscr C}\). Then \(R\) is a separated presheaf,
\(\ell_R:R\to LR\) is a covering monomorphism, and there is a commutative
diagram}
\[
\begin{array}{ccc}
 R&\xrightarrow{i}&F\\
 &\overset{\ell_R}{\searrow}
 &\overset{j}{\Big\uparrow}\\[-0.5ex]
 &&LR.
\end{array}
\]
\textit{Consequently, \(R\) is a refinement of \(F\) if and only if
\(j\) is an isomorphism.}

We already observed that \(R\) is separated and that \(j\) is a
monomorphism. By Proposition~\SGARefLink{sga3:iv:proposition:4:3:11:item:iv:01}{4.3.11(iv)},
\(\ell_R\) is a covering monomorphism. Thus, if \(j\) is an isomorphism,
\(R\) is a refinement of \(F\). Conversely, if \(i\) is covering, then so
is \(j\), and Proposition~\SGARefLink{sga3:iv:corollary:4:3:2:1}{4.3.2.1} shows that \(j\) is an isomorphism.

\medskip
\noindent\SGARefTarget{sga3:iv:remark:4:3:13}\textbf{Remark 4.3.13.}
If \(J'(S)\) is a cofinal subset of \(J(S)\), then
\SGARefTarget{sga3:iv:supp:remark:4:3:13:display:01}\[
 \operatorname{Hom}(S,LP)=
 \varinjlim_{R\in J'(S)}\operatorname{Hom}(R,P).
\]
In particular, let \(S\mapsto R(S)\) be a pretopology generating
the given topology. The functor \(L\) can be described using the covering
families belonging to \(R(S)\). Writing out \SGARefLink{sga3:iv:supp:remark:4:3:13:display:01}{the preceding formula}
recovers the construction of \SGARefLink{sga3:iv:bib:ma}{[MA]}.

Let \(i:\widetilde{\mathscr C}\to\widehat{\mathscr C}\) be the inclusion
functor. Proposition~\SGARefLink{sga3:iv:proposition:4:3:11}{4.3.11} gives \SGARefLink{sga3:iv:theorem:4:3:14}{the following theorem}.

\medskip
\noindent\SGARefTarget{sga3:iv:theorem:4:3:14}\textbf{Theorem 4.3.14.}
\textit{There is a unique functor
\(a:\widehat{\mathscr C}\to\widetilde{\mathscr C}\) for which the diagram}
\[
\begin{array}{ccc}
 \widehat{\mathscr C}&\xrightarrow{L}&\widehat{\mathscr C}\\
 \overset{a}{\Big\downarrow}&&
 \overset{L}{\Big\downarrow}\\
 \widetilde{\mathscr C}&\xrightarrow{i}&\widehat{\mathscr C}
\end{array}
\]
\textit{is commutative; that is, \(L(L(P))\) is a sheaf for every
presheaf \(P\). The functors \(i\) and \(a\) are adjoint: for every
presheaf \(P\) and every sheaf \(F\), there is an isomorphism, functorial in
\(P\) and \(F\),}
\[
 \operatorname{Hom}_{\widehat{\mathscr C}}(P,i(F))
 \simeq
 \operatorname{Hom}_{\widetilde{\mathscr C}}(a(P),F),
\]
\textit{or simply}
\[
 \operatorname{Hom}(P,F)\simeq\operatorname{Hom}(a(P),F).
\]

\medskip
\noindent\SGARefTarget{sga3:iv:definition:4:3:15}\textbf{Definition 4.3.15.}
\textit{The sheaf \(a(P)\) is called the sheaf \emph{associated} to the
presheaf \(P\).}

\medskip
\noindent\SGARefTarget{sga3:iv:remark:4:3:16}\textbf{Remark 4.3.16.}
Since \(L\) is constructed using inverse limits and filtered direct limits,
it commutes with finite inverse limits.

Moreover, under the identification \(L(P\times P)=LP\times LP\), the
morphism \(\ell_{P\times P}\) is identified with
\(\ell_P\times\ell_P\). Thus, for example, if \(P\) is a presheaf of
groups, then \(LP\) is canonically a presheaf of groups and the canonical
morphism \(P\to LP\) is a morphism of groups. The same holds for the
functor \(a\). Consequently, if \(P\) is a presheaf of groups and \(F\) a
sheaf of groups, there is an isomorphism
\[
 \operatorname{Hom}_{\widehat{\mathscr C}\text{-gr.}}(P,i(F))
 \simeq
 \operatorname{Hom}_{\widetilde{\mathscr C}\text{-gr.}}(a(P),F).
\]
See \SGARefLink{sga3:iv:bib:d}{[D]} for further details.

\medskip
\noindent\SGARefTarget{sga3:iv:para:4:3:17}\textbf{4.3.17.}
Let \(\mathscr V\) be any category. A \emph{presheaf on \(\mathscr C\)
with values in \(\mathscr V\)} is a contravariant functor from
\(\mathscr C\) to \(\mathscr V\). Before defining sheaves with values in
\(\mathscr V\), we recall the definition of the inverse limit of a
functor. Let \(\mathscr R\) and \(\mathscr V\) be categories, and let
\[
 F:\mathscr R^{\circ}\longrightarrow\mathscr V
\]
be a contravariant functor. Write \(\varprojlim F\) for the object of
\(\widehat{\mathscr V}\) defined by
\[
 \begin{aligned}
 (\varprojlim F)(X)
 &=\operatorname{Hom}_{\widehat{\mathscr V}}
     (X,\varprojlim F)\\
 &=\varprojlim_{U\in\operatorname{Ob}\mathscr R}
     \operatorname{Hom}_{\mathscr V}(F(U),X)\\
 &=\operatorname{Hom}(c_X,F),
 \end{aligned}
\]
where \(X\) varies over the objects of \(\mathscr V\), where \(c_X\)
denotes the contravariant functor from \(\mathscr R\) to \(\mathscr V\)
that sends every object to \(X\) and every arrow to
\(\operatorname{id}_X\), and where the last \(\operatorname{Hom}\) is
taken in \(\operatorname{Hom}(\mathscr R^{\circ},\mathscr V)\). If
\(\mathscr R\) has a final object \(e_{\mathscr R}\), then
\(\varprojlim F=F(e_{\mathscr R})\). If \(\mathscr V=\mathbf{Set}\), the
functor \(\varprojlim\) is identified with \(\Gamma\).

Let \(S\) be an object of \(\mathscr C\) and \(R\) a sieve of \(S\). Let
\(\overline R\) be the full subcategory of \(\mathscr C_{/S}\) whose
object set is \(E(R)\), and let
\(i_R:\overline R\to\mathscr C_{/S}\) be the canonical functor. A
presheaf \(P\) on \(\mathscr C\) with values in \(\mathscr V\) restricts
to a functor
\[
 P_S:(\mathscr C_{/S})^{\circ}\longrightarrow\mathscr V.
\]
The functor \(i_R\) induces a morphism in \(\widehat{\mathscr V}\):
\[
 P(S)=P_S(S)=\varprojlim P_S
 \longrightarrow\varprojlim(P_S\circ i_R).
\]
We write
\[
 \varprojlim_RP=\varprojlim(P_S\circ i_R)
\]
for the latter object of \(\widehat{\mathscr V}\). In view of \SGARefLink{sga3:iv:remark:4:1:6}{4.1.6},
Definition~\SGARefLink{sga3:iv:definition:4:3:1}{4.3.1} generalizes as follows.

\medskip
\noindent\SGARefTarget{sga3:iv:definition:4:3:18}\textbf{Definition 4.3.18.}
\textit{A presheaf \(P\) on \(\mathscr C\) with values in \(\mathscr V\)
is \emph{separated} (respectively, a \emph{sheaf}) if, for every
\(S\in\operatorname{Ob}\mathscr C\) and every \(R\in J(S)\), the canonical
morphism in \(\widehat{\mathscr V}\)}
\[
 P(S)\longrightarrow\varprojlim_RP
\]
\textit{is a monomorphism (respectively, an isomorphism).}

When \(\mathscr V\) is the category \(\mathbf{Gr}\) of groups (or any
other category of sets endowed with algebraic structures defined by finite
inverse limits), the following notions are equivalent (cf. \SGARefLink{sga3:iv:bib:d}{[D]}): a
presheaf on \(\mathscr C\) with values in \(\mathbf{Gr}\) whose underlying
presheaf of sets is a sheaf, and a group object in the category of sheaves
of sets. With these identifications understood, we shall always regard
sheaves with values in a category of sets endowed with algebraic structures
defined by finite inverse limits as sheaves of sets equipped, in
\(\widetilde{\mathscr C}\), with the corresponding algebraic structure.

